\chapter{The deploy loop}
\label{ch:deploy}

\chapterbrief
  {Build on your laptop, run on the Pi, in about five seconds.}
  {Why the feedback loop's length decides how you work.}
  {\code{make deploy \&\& make run} as one reflex.}

\section{The loop is the thing}

How long it takes to see the effect of a change determines how you behave.

\begin{center}\small
\begin{tabular}{@{}L{0.16\linewidth}L{0.76\linewidth}@{}}
\toprule
\thead{Loop} & \thead{How you end up working} \\
\midrule
5 s   & Try it. Wrong? Change one thing, try again. You experiment. \\
40 s  & You batch up three changes to avoid the wait --- and when it fails you no longer know which one did it. \\
3 min & You stop running it. You read the code and \emph{reason}, which is far more error-prone than looking. \\
\bottomrule
\end{tabular}
\end{center}

The middle row is the dangerous one, because it feels productive.

\begin{analogy}{Tasting while you cook}
A cook tastes constantly --- a pinch of salt, taste, another pinch. Small
corrections, immediate feedback.

Now imagine each taste took three minutes. You would salt once, hopefully, and
find out at the table. That is what a slow build does to programming: it turns
tasting into guessing.
\end{analogy}

\section{Why not just build on the Pi?}

You could. \code{git push} from the laptop, \code{git pull} on the Pi, \code{mvn
compile}. It works, and on a Zero 2 W it takes roughly forty seconds.

The trick is that \textbf{Java bytecode is architecture-independent}. A
\code{.class} file compiled on your laptop --- x86 or Apple Silicon, does not
matter --- runs unchanged on the Pi's ARM chip. The JVM is the portability layer.

So: compile on the fast machine, copy the output to the slow one.

\begin{center}
\begin{tikzpicture}[font=\footnotesize\sffamily,
  b/.style={draw=brewernavy, rounded corners=2pt, minimum width=30mm,
            minimum height=8mm, inner sep=3pt}]
  \node[b] (l) {laptop: \code{mvn compile}};
  \node[b, right=20mm of l] (r) {Pi: \code{java -cp ...}};
  \draw[-{Stealth[length=4pt]}, brewergreen, line width=1pt] (l) --
     node[above, font=\scriptsize\sffamily, color=brewergreen] {rsync}
     node[below, font=\scriptsize\itshape, color=brewergrey] {only changed bytes} (r);
  \node[font=\scriptsize\itshape, color=brewergrey, below=1mm of l] {2 s};
  \node[font=\scriptsize\itshape, color=brewergrey, below=1mm of r] {instant};
\end{tikzpicture}
\end{center}

\section{One-time: the dependency jars}

Your code needs Pi4J and SLF4J at runtime. Those jars change only when
\code{pom.xml} changes, so copy them once:

\begin{shell}
# on your laptop, in the project
make deploy-full
\end{shell}

That builds, resolves the dependency jars into \code{target/dependency/}, and
copies both to the Pi.

\section{Every time after that}

\begin{shell}
make deploy
make run
\end{shell}

\begin{trace}
==> building locally
==> copying classes to martin@brewer:~/midnight-brewer
==> done
=== toolchain check ===
java.version    = 21.0.11
os.arch         = aarch64            <- note: you are on the Pi now
major version   = 21
  -> toolchain OK.
\end{trace}

That \code{os.arch = aarch64} is the proof: you compiled on your laptop and the
result ran on ARM.

To run something else:

\begin{shell}
make run MAIN=com.midnightbrewer.diag.SpiProbe
# or equivalently
./scripts/run.sh com.midnightbrewer.diag.SpiProbe
\end{shell}

\section{What the scripts actually do}

Worth reading once, so they are not magic.

\subsection*{\texttt{deploy.sh}}

\begin{shell}
mvn -q -DskipTests compile
rsync -a --delete target/classes/ "$PI:midnight-brewer/target/classes/"
\end{shell}

\begin{description}[leftmargin=0pt, style=nextline]
  \item[\code{-DskipTests}] tests already ran on your laptop; running them again
        on the Pi would be slow and pointless.
  \item[\code{rsync -a}] archive mode: recursive, preserves timestamps and
        permissions.
  \item[\code{-{}-delete}] remove files on the Pi that no longer exist locally.
        Without it, a class you renamed lingers over there --- and stale classes
        produce genuinely baffling behaviour.
  \item[the trailing slashes] \code{target/classes/} means \emph{the contents
        of}. Without the slash rsync would create
        \code{classes/classes/}. This bites everyone once.
\end{description}

\subsection*{\texttt{run.sh}}

\begin{shell}
ssh -t "$PI" "cd ~/midnight-brewer && java -cp 'target/classes:target/dependency/*' $MAIN"
\end{shell}

\code{-t} forces a terminal, so \textbf{Ctrl-C reaches the program on the Pi}
instead of only killing your local ssh client and leaving the program running
over there. Without it you accumulate orphaned processes holding the SPI device,
and the next run fails with a resource-busy error you will not understand.

\begin{trap}{Stale classes}
If you change a file and behaviour does not change, suspect stale output before
suspecting anything else.

\begin{shell}
mvn clean compile && make deploy-full
\end{shell}

The symptom is unmistakable once you know it: the code in front of you
\emph{cannot} produce the output you are seeing.
\end{trap}

\section{The rhythm}

\begin{shell}
# edit code in your editor
mvn test                 # fast, local, no hardware
make deploy && make run  # on the real thing
\end{shell}

Most of your loop is the middle line. The bottom line is for when you need
actual hardware --- and by Volume One, Chapter 1, you will have arranged for that
to be as rarely as possible.

\begin{checkpoint}{The loop works}
\begin{shell}
make deploy-full
make run
\end{shell}

You see the toolchain check, printed by the Pi, reporting \code{aarch64}.

Now change one line in \code{ToolchainCheck.java} --- add a
\code{System.out.println("hello from the loop");} --- and run:

\begin{shell}
make deploy && make run
\end{shell}

Your new line appears. Time it: it should be around five seconds. That number is
what the rest of the project will feel like.
\end{checkpoint}
